\section{Introduction}\label{sec:intro}
The standard cosmological model explains a wide range of observations through general relativity supplemented by cold dark matter (CDM) and a cosmological constant. Its empirical success is substantial, particularly for the cosmic microwave background (CMB), baryon acoustic oscillations (BAO), gravitational lensing, and large-scale structure \cite{Planck2020VI,DESI2024VI,DESIDR2II2025,Madhavacheril2024,DESY3Cosmology2022}. At galactic scales, the tight relation between baryonic distributions and observed accelerations has motivated the modified Newtonian dynamics (MOND) program and relativistic attempts to reproduce its phenomenology \cite{Milgrom1983,Bekenstein2004,LelliMcGaughSchombert2016,McGaughLelliSchombert2016}. A central theoretical challenge is to retain the empirical galactic regularities while obtaining a viable relativistic theory with stable perturbations and a cosmological component capable of replacing the dynamical role normally assigned to particle CDM. This question is especially timely because contemporary expansion-history and structure analyses continue to sharpen tests of departures from the baseline model \cite{DiValentino2021,Schoneberg2022,PerivolaropoulosSkara2022,DES2024SN,EuclidWeakLensing2024}.

Aether scalar--tensor (AeST) gravity provides a single-metric setting containing a metric, a unit timelike vector, and a shift-symmetric scalar \cite{SkordisZlosnik2021,SkordisZlosnik2022}. Subsequent studies have examined its nonlinear Hamiltonian structure, cosmological dynamics, quasistatic spherical solutions, galaxy-cluster configurations, and compact stars \cite{BatakiSkordisZlosnik2024,RosaZlosnik2024,VerwayenSkordisBoehm2024,DurakovicSkordis2024,ReyesSakstein2024,ReyesSakstein2025}. These results motivate a direct test of whether a carefully specified scalar--vector sector can act as an effective cold component without introducing particle dark matter.

Late-time acceleration is treated here through the auxiliary constraint sector of VCDM, a Type-II minimally modified gravity construction whose cosmological applications and perturbative degree-of-freedom structure have been studied independently \cite{DeFeliceMukohyamaPookkillath2021,DeFeliceMaedaMukohyamaPookkillath2022,GanzMartensMukohyamaNamba2023,AkarsuEtAl2024VCDM,AkarsuEtAl2026VCDM}. This sector changes the homogeneous gravitational dynamics without adding an additional propagating scalar mode. Combining it with AeST allows the particle-CDM density and bare cosmological constant to be fixed to zero while retaining independent control of the effective cold and acceleration sectors.

The new theoretical result of this work concerns the scalar kinetic structure. A free function separable between timelike and spacelike scalar invariants admits a radiation-era slow mode that is not strictly damped on the physical early-time solution. We show that the minimal nonseparable operator $\Y\Gamma(\Q)$ removes this defect when $\Gamma$ is fixed by the already existing conserved scalar current. The construction therefore introduces no new field, no new free function, and no additional continuous fit parameter. On FLRW, $\Y=0$, so the calibrated homogeneous scalar background is unchanged; in the screened static configuration, $\Q=\Q_0$ and $\Gamma(\Q_0)=0$, so the galactic function is unchanged.

We carry this construction through a full theory-to-data validation. We derive the perturbation equations and their map into CLASS \cite{Lesgourgues2011CLASS,Blas2011CLASS}, audit the combined nonlinear constraint algebra, canonically normalize the cubic scalar interactions, compare the nonseparable theory against its separable control at fixed background, and only then optimize the frozen numerical theory against Planck 2018, DESI DR2, and Pantheon+ likelihoods \cite{Brout2022,Scolnic2022,DESIDR2II2025}. A separate profile in an explicit particle-CDM density tests whether the data require such a component once the effective AeST cold density is allowed to readjust independently.

The claims of this article are deliberately limited to the classical relativistic theory, its linear cosmology, the tested static galactic limit, and the stated effective-field-theory domain. We do not infer a fundamental quantum ultraviolet completion from the classical cutoff analysis, and we do not interpret a raw non-nested likelihood difference as a Gaussian discovery significance.
